% Options for packages loaded elsewhere
\PassOptionsToPackage{unicode}{hyperref}
\PassOptionsToPackage{hyphens}{url}
\PassOptionsToPackage{dvipsnames,svgnames,x11names}{xcolor}
%
\documentclass[
]{scrartcl}

\usepackage{amsmath,amssymb}
\usepackage{iftex}
\ifPDFTeX
  \usepackage[T1]{fontenc}
  \usepackage[utf8]{inputenc}
  \usepackage{textcomp} % provide euro and other symbols
\else % if luatex or xetex
  \usepackage{unicode-math}
  \defaultfontfeatures{Scale=MatchLowercase}
  \defaultfontfeatures[\rmfamily]{Ligatures=TeX,Scale=1}
\fi
\usepackage{lmodern}
\ifPDFTeX\else  
    % xetex/luatex font selection
\fi
% Use upquote if available, for straight quotes in verbatim environments
\IfFileExists{upquote.sty}{\usepackage{upquote}}{}
\IfFileExists{microtype.sty}{% use microtype if available
  \usepackage[]{microtype}
  \UseMicrotypeSet[protrusion]{basicmath} % disable protrusion for tt fonts
}{}
\makeatletter
\@ifundefined{KOMAClassName}{% if non-KOMA class
  \IfFileExists{parskip.sty}{%
    \usepackage{parskip}
  }{% else
    \setlength{\parindent}{0pt}
    \setlength{\parskip}{6pt plus 2pt minus 1pt}}
}{% if KOMA class
  \KOMAoptions{parskip=half}}
\makeatother
\usepackage{xcolor}
\setlength{\emergencystretch}{3em} % prevent overfull lines
\setcounter{secnumdepth}{5}
% Make \paragraph and \subparagraph free-standing
\ifx\paragraph\undefined\else
  \let\oldparagraph\paragraph
  \renewcommand{\paragraph}[1]{\oldparagraph{#1}\mbox{}}
\fi
\ifx\subparagraph\undefined\else
  \let\oldsubparagraph\subparagraph
  \renewcommand{\subparagraph}[1]{\oldsubparagraph{#1}\mbox{}}
\fi

% Lower the position of page numbers
\usepackage{fancyhdr}
\fancyhf{}
\fancyfoot[C]{\thepage}
\setlength{\footskip}{100pt} % Increase this value to lower the page numbers


\usepackage{color}
\usepackage{fancyvrb}
\newcommand{\VerbBar}{|}
\newcommand{\VERB}{\Verb[commandchars=\\\{\}]}
\DefineVerbatimEnvironment{Highlighting}{Verbatim}{commandchars=\\\{\}}
% Add ',fontsize=\small' for more characters per line
\usepackage{framed}
\definecolor{shadecolor}{RGB}{241,243,245}
\newenvironment{Shaded}{\begin{snugshade}}{\end{snugshade}}
\newcommand{\AlertTok}[1]{\textcolor[rgb]{0.68,0.00,0.00}{#1}}
\newcommand{\AnnotationTok}[1]{\textcolor[rgb]{0.37,0.37,0.37}{#1}}
\newcommand{\AttributeTok}[1]{\textcolor[rgb]{0.40,0.45,0.13}{#1}}
\newcommand{\BaseNTok}[1]{\textcolor[rgb]{0.68,0.00,0.00}{#1}}
\newcommand{\BuiltInTok}[1]{\textcolor[rgb]{0.00,0.23,0.31}{#1}}
\newcommand{\CharTok}[1]{\textcolor[rgb]{0.13,0.47,0.30}{#1}}
\newcommand{\CommentTok}[1]{\textcolor[rgb]{0.37,0.37,0.37}{#1}}
\newcommand{\CommentVarTok}[1]{\textcolor[rgb]{0.37,0.37,0.37}{\textit{#1}}}
\newcommand{\ConstantTok}[1]{\textcolor[rgb]{0.56,0.35,0.01}{#1}}
\newcommand{\ControlFlowTok}[1]{\textcolor[rgb]{0.00,0.23,0.31}{\textbf{#1}}}
\newcommand{\DataTypeTok}[1]{\textcolor[rgb]{0.68,0.00,0.00}{#1}}
\newcommand{\DecValTok}[1]{\textcolor[rgb]{0.68,0.00,0.00}{#1}}
\newcommand{\DocumentationTok}[1]{\textcolor[rgb]{0.37,0.37,0.37}{\textit{#1}}}
\newcommand{\ErrorTok}[1]{\textcolor[rgb]{0.68,0.00,0.00}{#1}}
\newcommand{\ExtensionTok}[1]{\textcolor[rgb]{0.00,0.23,0.31}{#1}}
\newcommand{\FloatTok}[1]{\textcolor[rgb]{0.68,0.00,0.00}{#1}}
\newcommand{\FunctionTok}[1]{\textcolor[rgb]{0.28,0.35,0.67}{#1}}
\newcommand{\ImportTok}[1]{\textcolor[rgb]{0.00,0.46,0.62}{#1}}
\newcommand{\InformationTok}[1]{\textcolor[rgb]{0.37,0.37,0.37}{#1}}
\newcommand{\KeywordTok}[1]{\textcolor[rgb]{0.00,0.23,0.31}{\textbf{#1}}}
\newcommand{\NormalTok}[1]{\textcolor[rgb]{0.00,0.23,0.31}{#1}}
\newcommand{\OperatorTok}[1]{\textcolor[rgb]{0.37,0.37,0.37}{#1}}
\newcommand{\OtherTok}[1]{\textcolor[rgb]{0.00,0.23,0.31}{#1}}
\newcommand{\PreprocessorTok}[1]{\textcolor[rgb]{0.68,0.00,0.00}{#1}}
\newcommand{\RegionMarkerTok}[1]{\textcolor[rgb]{0.00,0.23,0.31}{#1}}
\newcommand{\SpecialCharTok}[1]{\textcolor[rgb]{0.37,0.37,0.37}{#1}}
\newcommand{\SpecialStringTok}[1]{\textcolor[rgb]{0.13,0.47,0.30}{#1}}
\newcommand{\StringTok}[1]{\textcolor[rgb]{0.13,0.47,0.30}{#1}}
\newcommand{\VariableTok}[1]{\textcolor[rgb]{0.07,0.07,0.07}{#1}}
\newcommand{\VerbatimStringTok}[1]{\textcolor[rgb]{0.13,0.47,0.30}{#1}}
\newcommand{\WarningTok}[1]{\textcolor[rgb]{0.37,0.37,0.37}{\textit{#1}}}

\providecommand{\tightlist}{%
  \setlength{\itemsep}{0pt}\setlength{\parskip}{0pt}}\usepackage{longtable,booktabs,array}
\usepackage{calc} % for calculating minipage widths
% Correct order of tables after \paragraph or \subparagraph
\usepackage{etoolbox}
\makeatletter
\patchcmd\longtable{\par}{\if@noskipsec\mbox{}\fi\par}{}{}
\makeatother
% Allow footnotes in longtable head/foot
\IfFileExists{footnotehyper.sty}{\usepackage{footnotehyper}}{\usepackage{footnote}}
\makesavenoteenv{longtable}
\usepackage{graphicx}
\makeatletter
\newsavebox\pandoc@box
\newcommand*\pandocbounded[1]{% scales image to fit in text height/width
  \sbox\pandoc@box{#1}%
  \Gscale@div\@tempa{\textheight}{\dimexpr\ht\pandoc@box+\dp\pandoc@box\relax}%
  \Gscale@div\@tempb{\linewidth}{\wd\pandoc@box}%
  \ifdim\@tempb\p@<\@tempa\p@\let\@tempa\@tempb\fi% select the smaller of both
  \ifdim\@tempa\p@<\p@\scalebox{\@tempa}{\usebox\pandoc@box}%
  \else\usebox{\pandoc@box}%
  \fi%
}
% Set default figure placement to htbp
\def\fps@figure{htbp}
\makeatother

% Change font size
\usepackage[font={footnotesize, sf}, labelfont=bf, margin=0.05\linewidth]{caption}
\renewcommand{\familydefault}{\sfdefault}
% Bold the 'Figure #' in the caption and separate it from the title/caption
% with a period
% Captions will be left justified
\usepackage[
	aboveskip=1pt,
	labelfont=bf,
	labelsep=period,
	justification=raggedright,
	singlelinecheck=off
]{caption}
\renewcommand{\figurename}{Fig}

% Add numbered lines
\usepackage{lineno}
% \linenumbers

% Package to include multiple title pages
% This will allow me to add a tile to the main text and to the SI
\usepackage{titling}

% Define command to begin the supplementary section
\newcommand{\beginsupplement}{
\setcounter{page}{1} % Restart page counter
\setcounter{section}{0} % Restart section counter
% \renewcommand{\thesection}{\Alph{section}}%
\renewcommand{\thesection}{}
\renewcommand{\thesubsection}{\Alph{subsection}}
\setcounter{table}{0} % Restart table counter
\renewcommand{\thetable}{\Alph{table}}
\setcounter{figure}{0} % Restart figure counter
\renewcommand{\thefigure}{\Alph{figure}}%
\setcounter{equation}{0} % Restart equation counter
\renewcommand{\theequation}{S\arabic{equation}}%
}

% Add author affiliations
\usepackage{authblk}

% Allow to define personalized colors
\usepackage{xcolor}
\makeatletter
\@ifpackageloaded{caption}{}{\usepackage{caption}}
\AtBeginDocument{%
\ifdefined\contentsname
  \renewcommand*\contentsname{Table of contents}
\else
  \newcommand\contentsname{Table of contents}
\fi
\ifdefined\listfigurename
  \renewcommand*\listfigurename{List of Figures}
\else
  \newcommand\listfigurename{List of Figures}
\fi
\ifdefined\listtablename
  \renewcommand*\listtablename{List of Tables}
\else
  \newcommand\listtablename{List of Tables}
\fi
\ifdefined\figurename
  \renewcommand*\figurename{Figure}
\else
  \newcommand\figurename{Figure}
\fi
\ifdefined\tablename
  \renewcommand*\tablename{Table}
\else
  \newcommand\tablename{Table}
\fi
}
\@ifpackageloaded{float}{}{\usepackage{float}}
\floatstyle{ruled}
\@ifundefined{c@chapter}{\newfloat{codelisting}{h}{lop}}{\newfloat{codelisting}{h}{lop}[chapter]}
\floatname{codelisting}{Listing}
\newcommand*\listoflistings{\listof{codelisting}{List of Listings}}
\makeatother
\makeatletter
\makeatother
\makeatletter
\@ifpackageloaded{caption}{}{\usepackage{caption}}
\@ifpackageloaded{subcaption}{}{\usepackage{subcaption}}
\makeatother
\ifLuaTeX
  \usepackage{selnolig}  % disable illegal ligatures
\fi

%%%%%%%%%%%%%%%%%%%%%%%%%%%%%%%%%%%%%%%%%%%%%%%%%%%%%%%%%%%%%%%%%%%%%%%%%%%%%%%%
\usepackage[
    style=vancouver,
    sorting=none,				% Do not sort bibliography
	url=false, 					% Do not show url in reference
	doi=false, 					% Do not show doi in reference
	isbn=false, 				% Do not show isbn link in reference
	eprint=false, 			    % Do not show eprint link in reference
	maxbibnames=9, 			    % Include up to 9 names in citation
	firstinits=true,
]{biblatex}
%%%%%%%%%%%%%%%%%%%%%%%%%%%%%%%%%%%%%%%%%%%%%%%%%%%%%%%%%%%%%%%%%%%%%%%%%%%%%%%%
\addbibresource{../references.bib}
\IfFileExists{bookmark.sty}{\usepackage{bookmark}}{\usepackage{hyperref}}
\IfFileExists{xurl.sty}{\usepackage{xurl}}{} % add URL line breaks if available
\urlstyle{same} % disable monospaced font for URLs
\hypersetup{
  pdftitle={Training an RHVAE on synthetic data},
  pdfauthor={Manuel Razo-Mejia},
  pdfkeywords={variational autoencoders, Riemannian geometry, deep
learning},
  colorlinks=true,
  linkcolor={blue},
  filecolor={Maroon},
  citecolor={Blue},
  urlcolor={Blue},
  pdfcreator={LaTeX via pandoc}}

%%%%%%%%%%%%%%%%%%%%%%%%%%%%%%%%%%%%%%%%%%%%%%%%%%%%%%%%%%%%%%%%%%%%%%%%%%%%%%%%
% Add title
\title{Training an RHVAE on synthetic data}
%%%%%%%%%%%%%%%%%%%%%%%%%%%%%%%%%%%%%%%%%%%%%%%%%%%%%%%%%%%%%%%%%%%%%%%%%%%%%%%%

%%%%%%%%%%%%%%%%%%%%%%%%%%%%%%%%%%%%%%%%%%%%%%%%%%%%%%%%%%%%%%%%%%%%%%%%%%%%%%%%
% Add list of authors

% Loop through authors
\author[1]{Manuel Razo-Mejia}

% Loop through affiliations
\affil[1]{Department of Biology, Stanford University}

% Add corresponding author
%%%%%%%%%%%%%%%%%%%%%%%%%%%%%%%%%%%%%%%%%%%%%%%%%%%%%%%%%%%%%%%%%%%%%%%%%%%%%%%%

%%%%%%%%%%%%%%%%%%%%%%%%%%%%%%%%%%%%%%%%%%%%%%%%%%%%%%%%%%%%%%%%%%%%%%%%%%%%%%%%
% Set affiliations in small font
\renewcommand\Affilfont{\itshape\small}
%%%%%%%%%%%%%%%%%%%%%%%%%%%%%%%%%%%%%%%%%%%%%%%%%%%%%%%%%%%%%%%%%%%%%%%%%%%%%%%%

%%%%%%%%%%%%%%%%%%%%%%%%%%%%%%%%%%%%%%%%%%%%%%%%%%%%%%%%%%%%%%%%%%%%%%%%%%%%%%%%
% Remove date
\date{}
%%%%%%%%%%%%%%%%%%%%%%%%%%%%%%%%%%%%%%%%%%%%%%%%%%%%%%%%%%%%%%%%%%%%%%%%%%%%%%%%

% Add package for line numbering
\usepackage{lineno}

\begin{document}
\linenumbers
\maketitle
(c) This work is licensed under a
\href{https://creativecommons.org/licenses/by/4.0/}{Creative Commons
Attribution License CC-BY 4.0}. All code contained herein is licensed
under an \href{https://opensource.org/licenses/MIT}{MIT license}.

\section{Riemannian Hamiltonian Variational Autoencoder for Phenotypic
Space
Reconstruction}\label{riemannian-hamiltonian-variational-autoencoder-for-phenotypic-space-reconstruction}

This notebook implements a Riemannian Hamiltonian Variational
Autoencoder (RHVAE) to analyze fitness profiles generated from
evolutionary simulations, as the ones generated in the
\href{evolutionary_dynamics.qmd}{Evolutionary Dynamics} notebook. We'll
train the RHVAE to learn a low-dimensional representation (latent space)
of these fitness profiles, with the goal of reconstructing the
underlying phenotypic coordinates.

\subsection{Introduction to RHVAEs}\label{introduction-to-rhvaes}

Variational Autoencoders (VAEs) are generative models that learn to
encode high-dimensional data into a lower-dimensional latent space and
then decode it back. A standard VAE consists of an encoder network that
maps input data to a distribution in latent space, and a decoder network
that reconstructs the input from samples from this distribution.

The Riemannian Hamiltonian VAE (RHVAE) extends this framework by
modeling the latent space as a Riemannian manifold, learning not only
the encoding/decoding functions but also the geometric structure of the
latent space. This is achieved by:

\begin{enumerate}
\def\labelenumi{\arabic{enumi}.}
\tightlist
\item
  Learning a position-dependent metric tensor that captures how
  distances in latent space relate to distances in data space.
\item
  Using Hamiltonian dynamics to move in latent space in a geometry-aware
  manner.
\end{enumerate}

The RHVAE is particularly useful for our application because: - It
captures nonlinear relationships between fitness profiles and underlying
phenotypes. - It provides meaningful distances in latent space via the
learned metric tensor. - It offers a principled way to interpolate
between points in latent space.

\subsection{Setup Environment}\label{setup-environment}

First, let's import the necessary packages for our implementation.

\begin{Shaded}
\begin{Highlighting}[]
\CommentTok{\# Import project package}
\ImportTok{import} \BuiltInTok{Antibiotic}
\ImportTok{import} \BuiltInTok{Antibiotic.mh}\NormalTok{ as mh }\CommentTok{\# Metropolis{-}Hastings dynamics module}

\CommentTok{\# Import AutoEncoderToolkit to train VAEs}
\ImportTok{import} \BuiltInTok{AutoEncoderToolkit}\NormalTok{ as AET}

\CommentTok{\# Import packages for manipulating results}
\ImportTok{import} \BuiltInTok{DimensionalData}\NormalTok{ as DD}
\ImportTok{import} \BuiltInTok{DataFrames}\NormalTok{ as DF}
\ImportTok{import} \BuiltInTok{Glob}

\CommentTok{\# Import ML libraries}
\ImportTok{import} \BuiltInTok{Flux}

\CommentTok{\# Import CUDA (if available) to train using GPU}
\CommentTok{\# import CUDA}

\CommentTok{\# Import library to save models}
\ImportTok{import} \BuiltInTok{JLD2}

\CommentTok{\# Import libraries for data handling}
\ImportTok{import} \BuiltInTok{IterTools}
\ImportTok{import} \BuiltInTok{StatsBase}
\ImportTok{import} \BuiltInTok{Random}

\CommentTok{\# Load Plotting packages}
\ImportTok{using} \BuiltInTok{CairoMakie}
\ImportTok{import} \BuiltInTok{ColorSchemes}
\ImportTok{import} \BuiltInTok{Colors}

\CommentTok{\# Activate backend}
\NormalTok{CairoMakie.}\FunctionTok{activate!}\NormalTok{()}

\CommentTok{\# Set plotting style}
\NormalTok{Antibiotic.viz.}\FunctionTok{theme\_makie!}\NormalTok{()}

\CommentTok{\# Set random seed for reproducibility}
\BuiltInTok{Random}\NormalTok{.}\FunctionTok{seed!}\NormalTok{(}\FloatTok{42}\NormalTok{)}
\end{Highlighting}
\end{Shaded}

\subsection{Defining Directories}\label{defining-directories}

We'll set up the directories where we'll load data from and save our
model and results to.

\begin{Shaded}
\begin{Highlighting}[]
\CommentTok{\# Defining directories...}

\CommentTok{\# Define output directory}
\NormalTok{out\_dir }\OperatorTok{=} \StringTok{"./sim\_metropolis\_dynamics"}

\CommentTok{\# Define model state directory}
\NormalTok{state\_dir }\OperatorTok{=} \StringTok{"}\SpecialCharTok{$}\NormalTok{(out\_dir)}\StringTok{/rhvae\_model\_state"}
\end{Highlighting}
\end{Shaded}

\subsection{Model Architecture and
Hyperparameters}\label{model-architecture-and-hyperparameters}

The RHVAE architecture consists of three main components: 1.
\textbf{Encoder}: Maps fitness profiles to a distribution in latent
space 2. \textbf{Decoder}: Reconstructs fitness profiles from latent
space coordinates 3. \textbf{Metric Network}: Learns the Riemannian
metric tensor of the latent space

Let's define the hyperparameters and architecture of our model

\begin{Shaded}
\begin{Highlighting}[]
\CommentTok{\# Defining hyperparameters...}

\CommentTok{\# Define dimensionality of latent space}
\NormalTok{n\_latent }\OperatorTok{=} \FloatTok{2}
\CommentTok{\# Define number of neurons in hidden layers}
\NormalTok{n\_neuron }\OperatorTok{=} \FloatTok{128}

\CommentTok{\# Define RHVAE hyper{-}parameters}
\NormalTok{T }\OperatorTok{=} \FloatTok{0.8f0} \CommentTok{\# Temperature}
\NormalTok{λ }\OperatorTok{=} \FloatTok{1.0f{-}2} \CommentTok{\# Regularization parameter}
\NormalTok{n\_centroids }\OperatorTok{=} \FloatTok{256} \CommentTok{\# Number of centroids}

\CommentTok{\# Define loss function hyper{-}parameters}
\NormalTok{ϵ }\OperatorTok{=} \FunctionTok{Float32}\NormalTok{(}\FloatTok{1E{-}3}\NormalTok{) }\CommentTok{\# Leapfrog step size}
\NormalTok{K }\OperatorTok{=} \FloatTok{10} \CommentTok{\# Number of leapfrog steps}
\NormalTok{βₒ }\OperatorTok{=} \FloatTok{0.3f0} \CommentTok{\# Initial temperature for tempering}

\CommentTok{\# Define RHVAE hyper{-}parameters in a named tuple}
\NormalTok{rhvae\_kwargs }\OperatorTok{=}\NormalTok{ (K}\OperatorTok{=}\NormalTok{K, ϵ}\OperatorTok{=}\NormalTok{ϵ, βₒ}\OperatorTok{=}\NormalTok{βₒ)}

\CommentTok{\# Define training hyperparameters}
\NormalTok{n\_epoch }\OperatorTok{=} \FloatTok{75} \CommentTok{\# Number of epochs}
\NormalTok{n\_batch }\OperatorTok{=} \FloatTok{512} \CommentTok{\# Batch size}
\NormalTok{n\_batch\_loss }\OperatorTok{=} \FloatTok{512} \CommentTok{\# Batch size for loss computation}
\NormalTok{η }\OperatorTok{=} \FloatTok{10}\OperatorTok{\^{}{-}}\FloatTok{3} \CommentTok{\# Learning rate}
\NormalTok{split\_frac }\OperatorTok{=} \FloatTok{0.85} \CommentTok{\# Train/val split}

\CommentTok{\# Define ELBO prefactors}
\NormalTok{logp\_prefactor }\OperatorTok{=}\NormalTok{ [}\FloatTok{10.0f0}\NormalTok{, }\FloatTok{0.1f0}\NormalTok{, }\FloatTok{0.1f0}\NormalTok{]}
\NormalTok{logq\_prefactor }\OperatorTok{=}\NormalTok{ [}\FloatTok{0.1f0}\NormalTok{, }\FloatTok{0.1f0}\NormalTok{, }\FloatTok{0.1f0}\NormalTok{]}

\CommentTok{\# Define loss function kwargs in a NamedTuple}
\NormalTok{loss\_kwargs }\OperatorTok{=}\NormalTok{ (}
\NormalTok{    K}\OperatorTok{=}\NormalTok{K,}
\NormalTok{    ϵ}\OperatorTok{=}\NormalTok{ϵ,}
\NormalTok{    βₒ}\OperatorTok{=}\NormalTok{βₒ,}
\NormalTok{    logp\_prefactor}\OperatorTok{=}\NormalTok{logp\_prefactor,}
\NormalTok{    logq\_prefactor}\OperatorTok{=}\NormalTok{logq\_prefactor,}
\NormalTok{)}

\CommentTok{\# Define by how much to subsample the time series}
\NormalTok{n\_sub }\OperatorTok{=} \FloatTok{10}
\end{Highlighting}
\end{Shaded}

\subsection{Loading and Preprocessing
Data}\label{loading-and-preprocessing-data}

Now let's load the fitness profiles from our evolutionary simulations
and preprocess them for training. Since the simulations are very densly
sampled over time, we will train the RHVAE on a subsampled time series,
taking every \texttt{n\_sub} time point.

\begin{Shaded}
\begin{Highlighting}[]
\CommentTok{\# Loading data into memory...}

\CommentTok{\# Load fitnotype profiles}
\NormalTok{fitnotype\_profiles }\OperatorTok{=}\NormalTok{ JLD2.}\FunctionTok{load}\NormalTok{(}\StringTok{"}\SpecialCharTok{$}\NormalTok{(out\_dir)}\StringTok{/sim\_evo.jld2"}\NormalTok{)[}\StringTok{"fitnotype\_profiles"}\NormalTok{]}

\CommentTok{\# Extract initial and final time points}
\NormalTok{t\_init, t\_final }\OperatorTok{=} \FunctionTok{collect}\NormalTok{(DD.}\FunctionTok{dims}\NormalTok{(fitnotype\_profiles, }\OperatorTok{:}\NormalTok{time)[[}\FloatTok{1}\NormalTok{, }\KeywordTok{end}\NormalTok{]])}
\CommentTok{\# Subsample time series}
\NormalTok{fitnotype\_profiles }\OperatorTok{=}\NormalTok{ fitnotype\_profiles[time}\OperatorTok{=}\NormalTok{DD.}\FunctionTok{At}\NormalTok{(t\_init}\OperatorTok{:}\NormalTok{n\_sub}\OperatorTok{:}\NormalTok{t\_final)]}

\CommentTok{\# Define number of environments}
\NormalTok{n\_env }\OperatorTok{=} \FunctionTok{length}\NormalTok{(DD.}\FunctionTok{dims}\NormalTok{(fitnotype\_profiles, }\OperatorTok{:}\NormalTok{landscape))}

\CommentTok{\# Extract fitness data bringing the fitness dimension to the first dimension}
\NormalTok{fit\_data }\OperatorTok{=} \FunctionTok{permutedims}\NormalTok{(fitnotype\_profiles.fitness.data, (}\FloatTok{5}\NormalTok{, }\FloatTok{1}\NormalTok{, }\FloatTok{2}\NormalTok{, }\FloatTok{3}\NormalTok{, }\FloatTok{4}\NormalTok{, }\FloatTok{6}\NormalTok{))}
\CommentTok{\# Reshape the array to a Matrix}
\NormalTok{fit\_data }\OperatorTok{=} \FunctionTok{reshape}\NormalTok{(fit\_data, }\FunctionTok{size}\NormalTok{(fit\_data, }\FloatTok{1}\NormalTok{), }\OperatorTok{:}\NormalTok{)}

\CommentTok{\# Reshape the array to stack the 3rd dimension}
\NormalTok{fit\_mat }\OperatorTok{=} \FunctionTok{log}\NormalTok{.(fit\_data)}

\CommentTok{\# Fit model to standardize data to mean zero and standard deviation 1 on each}
\CommentTok{\# environment }
\NormalTok{dt }\OperatorTok{=}\NormalTok{ StatsBase.}\FunctionTok{fit}\NormalTok{(StatsBase.ZScoreTransform, fit\_mat, dims}\OperatorTok{=}\FloatTok{2}\NormalTok{)}

\CommentTok{\# Standardize the data to have mean 0 and standard deviation 1}
\NormalTok{fit\_std }\OperatorTok{=}\NormalTok{ StatsBase.}\FunctionTok{transform}\NormalTok{(dt, fit\_mat)}

\CommentTok{\# Split indexes of data into training and validation}
\NormalTok{train\_idx, val\_idx }\OperatorTok{=}\NormalTok{ Flux.}\FunctionTok{splitobs}\NormalTok{(}
    \FloatTok{1}\OperatorTok{:}\FunctionTok{size}\NormalTok{(fit\_std, }\FloatTok{2}\NormalTok{), at}\OperatorTok{=}\NormalTok{split\_frac, shuffle}\OperatorTok{=}\ConstantTok{true}
\NormalTok{)}

\CommentTok{\# Extract train and validation data}
\NormalTok{train\_data }\OperatorTok{=}\NormalTok{ fit\_std[}\OperatorTok{:}\NormalTok{, train\_idx]}
\NormalTok{val\_data }\OperatorTok{=}\NormalTok{ fit\_std[}\OperatorTok{:}\NormalTok{, val\_idx]}
\end{Highlighting}
\end{Shaded}

Our data preprocessing involves several key steps: 1. \textbf{Log
transformation}: We apply a logarithmic transformation to the fitness
values to make their distribution more amenable to the model. 2.
\textbf{Standardization}: We standardize the log-fitness values per
environment to have mean 0 and standard deviation 1. This is common
practice when training neural networks as it helps with numerical
stability and convergence. 3. \textbf{Train-validation split}: We split
the data into training (85\%) and validation (15\%) sets. The validation
set is used to monitor the model's performance during training.

\subsection{Defining the RHVAE Model
Architecture}\label{defining-the-rhvae-model-architecture}

We now define the RHVAE model architecture with its three core
components: encoder, decoder, and metric network. For this, we use the
\texttt{AutoEncoderToolkit.jl} package \autocite{razo-mejia2024}.

The first step is to select the centroids used to ``anchor'' the metric
tensor learned by the RHVAE. We do this by using the
\texttt{centroids\_kmedoids} function from the
\texttt{AutoEncoderToolkit.jl} package. These centroids will be part of
the definition of the RHVAE model itself.

\begin{Shaded}
\begin{Highlighting}[]
\CommentTok{\# Selecting centroids via k{-}means...}

\CommentTok{\# Select centroids via k{-}medoids}
\NormalTok{centroids\_data }\OperatorTok{=}\NormalTok{ AET.utils.}\FunctionTok{centroids\_kmedoids}\NormalTok{(fit\_std, n\_centroids)}
\end{Highlighting}
\end{Shaded}

Now, we can define the RHVAE architecture. We start by defining the
encoder, which maps the fitness profiles to a distribution in latent
space. Since \texttt{AutoEncoderToolkit.jl} is based on
\texttt{Flux.jl}, we will build all of the components of the RHVAE using
\texttt{Flux.jl} components.

\begin{Shaded}
\begin{Highlighting}[]
\CommentTok{\# Define JointGaussianLogEncoder...}

\CommentTok{\# Define encoder chain}
\NormalTok{encoder\_chain }\OperatorTok{=}\NormalTok{ Flux.}\FunctionTok{Chain}\NormalTok{(}
    \CommentTok{\# First layer}
\NormalTok{    Flux.}\FunctionTok{Dense}\NormalTok{(n\_env }\OperatorTok{=\textgreater{}}\NormalTok{ n\_neuron, Flux.identity),}
    \CommentTok{\# Second layer}
\NormalTok{    Flux.}\FunctionTok{Dense}\NormalTok{(n\_neuron }\OperatorTok{=\textgreater{}}\NormalTok{ n\_neuron, Flux.leakyrelu),}
    \CommentTok{\# Third layer}
\NormalTok{    Flux.}\FunctionTok{Dense}\NormalTok{(n\_neuron }\OperatorTok{=\textgreater{}}\NormalTok{ n\_neuron, Flux.leakyrelu),}
    \CommentTok{\# Fourth layer}
\NormalTok{    Flux.}\FunctionTok{Dense}\NormalTok{(n\_neuron }\OperatorTok{=\textgreater{}}\NormalTok{ n\_neuron, Flux.leakyrelu),}
\NormalTok{)}

\CommentTok{\# Define layers for µ and log(σ)}
\NormalTok{µ\_layer }\OperatorTok{=}\NormalTok{ Flux.}\FunctionTok{Dense}\NormalTok{(n\_neuron }\OperatorTok{=\textgreater{}}\NormalTok{ n\_latent, Flux.identity)}
\NormalTok{logσ\_layer }\OperatorTok{=}\NormalTok{ Flux.}\FunctionTok{Dense}\NormalTok{(n\_neuron }\OperatorTok{=\textgreater{}}\NormalTok{ n\_latent, Flux.identity)}

\CommentTok{\# build encoder}
\NormalTok{encoder }\OperatorTok{=}\NormalTok{ AET.}\FunctionTok{JointGaussianLogEncoder}\NormalTok{(encoder\_chain, µ\_layer, logσ\_layer)}
\end{Highlighting}
\end{Shaded}

Now, we define the decoder, which maps the latent space coordinates to
the fitness profiles.

\begin{Shaded}
\begin{Highlighting}[]
\CommentTok{\# Define SimpleGaussianDecoder...}

\CommentTok{\# Initialize decoder}
\NormalTok{decoder }\OperatorTok{=}\NormalTok{ AET.}\FunctionTok{SimpleGaussianDecoder}\NormalTok{(}
\NormalTok{    Flux.}\FunctionTok{Chain}\NormalTok{(}
        \CommentTok{\# First layer}
\NormalTok{        Flux.}\FunctionTok{Dense}\NormalTok{(n\_latent }\OperatorTok{=\textgreater{}}\NormalTok{ n\_neuron, Flux.identity),}
        \CommentTok{\# Second Layer}
\NormalTok{        Flux.}\FunctionTok{Dense}\NormalTok{(n\_neuron }\OperatorTok{=\textgreater{}}\NormalTok{ n\_neuron, Flux.leakyrelu),}
        \CommentTok{\# Third layer}
\NormalTok{        Flux.}\FunctionTok{Dense}\NormalTok{(n\_neuron }\OperatorTok{=\textgreater{}}\NormalTok{ n\_neuron, Flux.leakyrelu),}
        \CommentTok{\# Fourth layer}
\NormalTok{        Flux.}\FunctionTok{Dense}\NormalTok{(n\_neuron }\OperatorTok{=\textgreater{}}\NormalTok{ n\_neuron, Flux.leakyrelu),}
        \CommentTok{\# Output layer}
\NormalTok{        Flux.}\FunctionTok{Dense}\NormalTok{(n\_neuron }\OperatorTok{=\textgreater{}}\NormalTok{ n\_env, Flux.identity)}
\NormalTok{    )}
\NormalTok{)}
\end{Highlighting}
\end{Shaded}

Finally, we define the metric network, which learns the Riemannian
metric tensor of the latent space.

\begin{Shaded}
\begin{Highlighting}[]
\CommentTok{\# Define MetricChain...}

\CommentTok{\# Define mlp chain}
\NormalTok{mlp\_chain }\OperatorTok{=}\NormalTok{ Flux.}\FunctionTok{Chain}\NormalTok{(}
    \CommentTok{\# First layer}
\NormalTok{    Flux.}\FunctionTok{Dense}\NormalTok{(n\_env }\OperatorTok{=\textgreater{}}\NormalTok{ n\_neuron, Flux.identity),}
    \CommentTok{\# Second layer}
\NormalTok{    Flux.}\FunctionTok{Dense}\NormalTok{(n\_neuron }\OperatorTok{=\textgreater{}}\NormalTok{ n\_neuron, Flux.leakyrelu),}
    \CommentTok{\# Third layer}
\NormalTok{    Flux.}\FunctionTok{Dense}\NormalTok{(n\_neuron }\OperatorTok{=\textgreater{}}\NormalTok{ n\_neuron, Flux.leakyrelu),}
    \CommentTok{\# Fourth layer}
\NormalTok{    Flux.}\FunctionTok{Dense}\NormalTok{(n\_neuron }\OperatorTok{=\textgreater{}}\NormalTok{ n\_neuron, Flux.leakyrelu),}
\NormalTok{)}

\CommentTok{\# Define layers for the diagonal and lower triangular part of the covariance}
\CommentTok{\# matrix}
\NormalTok{diag }\OperatorTok{=}\NormalTok{ Flux.}\FunctionTok{Dense}\NormalTok{(n\_neuron }\OperatorTok{=\textgreater{}}\NormalTok{ n\_latent, Flux.identity)}
\NormalTok{lower }\OperatorTok{=}\NormalTok{ Flux.}\FunctionTok{Dense}\NormalTok{(}
\NormalTok{    n\_neuron }\OperatorTok{=\textgreater{}}\NormalTok{ n\_latent }\OperatorTok{*}\NormalTok{ (n\_latent }\OperatorTok{{-}} \FloatTok{1}\NormalTok{) }\OperatorTok{÷} \FloatTok{2}\NormalTok{, Flux.identity}
\NormalTok{)}

\CommentTok{\# Build metric chain}
\NormalTok{metric\_chain }\OperatorTok{=}\NormalTok{ AET.RHVAEs.}\FunctionTok{MetricChain}\NormalTok{(mlp\_chain, diag, lower)}
\end{Highlighting}
\end{Shaded}

For the final step, we bring all the components together to define the
RHVAE model itself.

\begin{Shaded}
\begin{Highlighting}[]
\CommentTok{\# Initialize rhvae}
\NormalTok{rhvae }\OperatorTok{=}\NormalTok{ AET.RHVAEs.}\FunctionTok{RHVAE}\NormalTok{(}
\NormalTok{    encoder }\OperatorTok{*}\NormalTok{ decoder,}
\NormalTok{    metric\_chain,}
\NormalTok{    centroids\_data,}
\NormalTok{    T,}
\NormalTok{    λ}
\NormalTok{)}
\end{Highlighting}
\end{Shaded}

The RHVAE architecture comprises:

\begin{enumerate}
\def\labelenumi{\arabic{enumi}.}
\tightlist
\item
  \textbf{Encoder}:

  \begin{itemize}
  \tightlist
  \item
    Takes fitness profiles (dimension = number of environments) as input
  \item
    Processes through 4 hidden layers with 128 neurons each
  \item
    Outputs mean (μ) and log-variance (log σ) parameters of a diagonal
    Gaussian distribution in latent space
  \end{itemize}
\item
  \textbf{Decoder}:

  \begin{itemize}
  \tightlist
  \item
    Takes latent space coordinates (dimension = 2) as input
  \item
    Processes through 4 hidden layers with 128 neurons each
  \item
    Outputs reconstructed fitness profiles
  \end{itemize}
\item
  \textbf{Metric Network}:

  \begin{itemize}
  \tightlist
  \item
    Learns the parameters of the Riemannian metric tensor
  \item
    Takes decoded fitness profiles as input
  \item
    Outputs parameters for the diagonal and lower triangular parts of
    the metric tensor
  \end{itemize}
\end{enumerate}

We also select centroids using k-medoids clustering to approximate the
data manifold, which helps in computing the metric tensor more
efficiently.

\subsection{Training the RHVAE}\label{training-the-rhvae}

Now, let's train our RHVAE model on the fitness profiles:

\begin{Shaded}
\begin{Highlighting}[]
\FunctionTok{println}\NormalTok{(}\StringTok{"Checking previous model states..."}\NormalTok{)}

\CommentTok{\# List previous model parameters}
\NormalTok{model\_states }\OperatorTok{=} \FunctionTok{sort}\NormalTok{(Glob.}\FunctionTok{glob}\NormalTok{(}\StringTok{"}\SpecialCharTok{$}\NormalTok{(state\_dir)}\StringTok{/beta{-}rhvae\_epoch*.jld2"}\NormalTok{[}\FloatTok{2}\OperatorTok{:}\KeywordTok{end}\NormalTok{], }\StringTok{"/"}\NormalTok{))}

\CommentTok{\# Check if model states exist}
\ControlFlowTok{if} \FunctionTok{length}\NormalTok{(model\_states) }\OperatorTok{\textgreater{}} \FloatTok{0}
    \CommentTok{\# Load model state}
\NormalTok{    model\_state }\OperatorTok{=}\NormalTok{ JLD2.}\FunctionTok{load}\NormalTok{(model\_states[}\KeywordTok{end}\NormalTok{])[}\StringTok{"model\_state"}\NormalTok{]}
    \CommentTok{\# Input parameters to model}
\NormalTok{    Flux.}\FunctionTok{loadmodel!}\NormalTok{(rhvae, model\_state)}
    \CommentTok{\# Update metric parameters}
\NormalTok{    AET.RHVAEs.}\FunctionTok{update\_metric!}\NormalTok{(rhvae)}
    \CommentTok{\# Extract epoch number}
\NormalTok{    epoch\_init }\OperatorTok{=} \FunctionTok{parse}\NormalTok{(}
        \DataTypeTok{Int}\NormalTok{, }\FunctionTok{match}\NormalTok{(}\StringTok{r"epoch}\CharTok{(}\SpecialCharTok{\textbackslash{}d+}\CharTok{)}\StringTok{"}\NormalTok{, model\_states[}\KeywordTok{end}\NormalTok{]).captures[}\FloatTok{1}\NormalTok{]}
\NormalTok{    ) }\OperatorTok{+} \FloatTok{1}
\ControlFlowTok{else}
\NormalTok{    epoch\_init }\OperatorTok{=} \FloatTok{1}
\ControlFlowTok{end} \CommentTok{\# if}

\FunctionTok{println}\NormalTok{(}\StringTok{"Initial epoch: }\SpecialCharTok{$}\NormalTok{epoch\_init}\StringTok{"}\NormalTok{)}

\FunctionTok{println}\NormalTok{(}\StringTok{"Uploading model to GPU..."}\NormalTok{)}

\CommentTok{\# Check if CUDA is available}
\ControlFlowTok{if}\NormalTok{ CUDA.}\FunctionTok{functional}\NormalTok{()}
    \CommentTok{\# Upload model to GPU}
\NormalTok{    rhvae }\OperatorTok{=}\NormalTok{ Flux.}\FunctionTok{gpu}\NormalTok{(rhvae)}
    \CommentTok{\# Upload data to GPU}
\NormalTok{    train\_data }\OperatorTok{=}\NormalTok{ Flux.}\FunctionTok{gpu}\NormalTok{(train\_data)}
\NormalTok{    val\_data }\OperatorTok{=}\NormalTok{ Flux.}\FunctionTok{gpu}\NormalTok{(val\_data)}
\ControlFlowTok{end}

\CommentTok{\# Explicit setup of optimizer}
\NormalTok{opt\_rhvae }\OperatorTok{=}\NormalTok{ Flux.Train.}\FunctionTok{setup}\NormalTok{(}
\NormalTok{    Flux.Optimisers.}\FunctionTok{Adam}\NormalTok{(η),}
\NormalTok{    rhvae}
\NormalTok{)}

\FunctionTok{println}\NormalTok{(}\StringTok{"}\SpecialCharTok{\textbackslash{}n}\StringTok{Training RHVAE...}\SpecialCharTok{\textbackslash{}n}\StringTok{"}\NormalTok{)}

\CommentTok{\# Loop through number of epochs}
\ControlFlowTok{for}\NormalTok{ epoch }\KeywordTok{in}\NormalTok{ epoch\_init}\OperatorTok{:}\NormalTok{n\_epoch}
    \CommentTok{\# Define number of batches}
\NormalTok{    num\_batches }\OperatorTok{=} \FunctionTok{size}\NormalTok{(train\_data, }\FloatTok{2}\NormalTok{) }\OperatorTok{÷}\NormalTok{ n\_batch}
    \CommentTok{\# Shuffle data indexes}
\NormalTok{    idx\_shuffle }\OperatorTok{=} \BuiltInTok{Random}\NormalTok{.}\FunctionTok{shuffle}\NormalTok{(}\FloatTok{1}\OperatorTok{:}\FunctionTok{size}\NormalTok{(train\_data, }\FloatTok{2}\NormalTok{))}
    \CommentTok{\# Split indexes into batches}
\NormalTok{    idx\_batches }\OperatorTok{=}\NormalTok{ IterTools.}\FunctionTok{partition}\NormalTok{(idx\_shuffle, n\_batch)}
    \CommentTok{\# Loop through batches}
    \ControlFlowTok{for}\NormalTok{ (i, idx\_tuple) }\KeywordTok{in} \FunctionTok{enumerate}\NormalTok{(idx\_batches)}
        \FunctionTok{println}\NormalTok{(}\StringTok{"Epoch: }\SpecialCharTok{$}\NormalTok{(epoch)}\StringTok{ | Batch: }\SpecialCharTok{$}\NormalTok{(i)}\StringTok{ / }\SpecialCharTok{$}\NormalTok{(}\FunctionTok{length}\NormalTok{(idx\_batches))}\StringTok{"}\NormalTok{)}
        \CommentTok{\# Extract indexes}
\NormalTok{        idx\_batch }\OperatorTok{=} \FunctionTok{collect}\NormalTok{(idx\_tuple)}
        \CommentTok{\# Train RHVAE}
\NormalTok{        loss\_epoch }\OperatorTok{=}\NormalTok{ AET.RHVAEs.}\FunctionTok{train!}\NormalTok{(}
\NormalTok{            rhvae, train\_data[}\OperatorTok{:}\NormalTok{, idx\_batch], opt\_rhvae;}
\NormalTok{            loss\_kwargs}\OperatorTok{=}\NormalTok{loss\_kwargs, verbose}\OperatorTok{=}\ConstantTok{false}\NormalTok{, loss\_return}\OperatorTok{=}\ConstantTok{true}
\NormalTok{        )}
        \FunctionTok{println}\NormalTok{(}\StringTok{"Loss: }\SpecialCharTok{$}\NormalTok{(loss\_epoch)}\StringTok{"}\NormalTok{)}
    \ControlFlowTok{end} \CommentTok{\# for train\_loader}

    \CommentTok{\# Sample train data}
\NormalTok{    train\_sample }\OperatorTok{=}\NormalTok{ train\_data[}
        \OperatorTok{:}\NormalTok{,}
\NormalTok{        StatsBase.}\FunctionTok{sample}\NormalTok{(}\FloatTok{1}\OperatorTok{:}\FunctionTok{size}\NormalTok{(train\_data, }\FloatTok{2}\NormalTok{), n\_batch\_loss, replace}\OperatorTok{=}\ConstantTok{false}\NormalTok{)}
\NormalTok{    ]}
    \CommentTok{\# Sample val data}
\NormalTok{    val\_sample }\OperatorTok{=}\NormalTok{ val\_data}

    \FunctionTok{println}\NormalTok{(}\StringTok{"Computing loss in training and validation data..."}\NormalTok{)}
\NormalTok{    loss\_train }\OperatorTok{=}\NormalTok{ AET.RHVAEs.}\FunctionTok{loss}\NormalTok{(rhvae, train\_sample; loss\_kwargs}\OperatorTok{...}\NormalTok{)}
\NormalTok{    loss\_val }\OperatorTok{=}\NormalTok{ AET.RHVAEs.}\FunctionTok{loss}\NormalTok{(rhvae, val\_sample; loss\_kwargs}\OperatorTok{...}\NormalTok{)}

    \CommentTok{\# Forward pass sample through model}
    \FunctionTok{println}\NormalTok{(}\StringTok{"Computing MSE in training and validation data..."}\NormalTok{)}
\NormalTok{    out\_train }\OperatorTok{=} \FunctionTok{rhvae}\NormalTok{(train\_sample; rhvae\_kwargs}\OperatorTok{...}\NormalTok{).μ}
\NormalTok{    mse\_train }\OperatorTok{=}\NormalTok{ Flux.}\FunctionTok{mse}\NormalTok{(train\_sample, out\_train)}
\NormalTok{    out\_val }\OperatorTok{=} \FunctionTok{rhvae}\NormalTok{(val\_sample; rhvae\_kwargs}\OperatorTok{...}\NormalTok{).μ}
\NormalTok{    mse\_val }\OperatorTok{=}\NormalTok{ Flux.}\FunctionTok{mse}\NormalTok{(val\_sample, out\_val)}

    \FunctionTok{println}\NormalTok{(}
        \StringTok{"}\SpecialCharTok{\textbackslash{}n}\StringTok{ Epoch: }\SpecialCharTok{$}\NormalTok{(epoch)}\StringTok{ / }\SpecialCharTok{$}\NormalTok{(n\_epoch)}\SpecialCharTok{\textbackslash{}n}\StringTok{ "} \OperatorTok{*}
        \StringTok{"   {-} loss\_train: }\SpecialCharTok{$}\NormalTok{(loss\_train)}\SpecialCharTok{\textbackslash{}n}\StringTok{"} \OperatorTok{*}
        \StringTok{"   {-} loss\_val: }\SpecialCharTok{$}\NormalTok{(loss\_val)}\SpecialCharTok{\textbackslash{}n}\StringTok{"} \OperatorTok{*}
        \StringTok{"   {-} mse\_train: }\SpecialCharTok{$}\NormalTok{(mse\_train)}\SpecialCharTok{\textbackslash{}n}\StringTok{"} \OperatorTok{*}
        \StringTok{"   {-} mse\_val: }\SpecialCharTok{$}\NormalTok{(mse\_val)}\SpecialCharTok{\textbackslash{}n}\StringTok{"}
\NormalTok{    )}

    \CommentTok{\# Save checkpoint}
\NormalTok{    JLD2.}\FunctionTok{jldsave}\NormalTok{(}
        \StringTok{"}\SpecialCharTok{$}\NormalTok{(state\_dir)}\StringTok{/beta{-}rhvae\_epoch}\SpecialCharTok{$}\NormalTok{(}\FunctionTok{lpad}\NormalTok{(epoch, }\FloatTok{5}\NormalTok{, }\StringTok{"0"}\NormalTok{))}\StringTok{.jld2"}\NormalTok{,}
\NormalTok{        model\_state}\OperatorTok{=}\NormalTok{Flux.}\FunctionTok{state}\NormalTok{(rhvae) }\OperatorTok{|\textgreater{}}\NormalTok{ Flux.cpu,}
\NormalTok{        loss\_train}\OperatorTok{=}\NormalTok{loss\_train,}
\NormalTok{        loss\_val}\OperatorTok{=}\NormalTok{loss\_val,}
\NormalTok{        mse\_train}\OperatorTok{=}\NormalTok{mse\_train,}
\NormalTok{        mse\_val}\OperatorTok{=}\NormalTok{mse\_val,}
\NormalTok{        train\_idx}\OperatorTok{=}\NormalTok{train\_idx,}
\NormalTok{        val\_idx}\OperatorTok{=}\NormalTok{val\_idx,}
\NormalTok{    )}
\ControlFlowTok{end} \CommentTok{\# for n\_epoch}
\end{Highlighting}
\end{Shaded}

The training process involves:

\begin{enumerate}
\def\labelenumi{\arabic{enumi}.}
\tightlist
\item
  \textbf{Batched Training}: We process the data in batches of size 512.
\item
  \textbf{Loss Computation}: We use a custom loss function for RHVAEs
  that combines:

  \begin{itemize}
  \tightlist
  \item
    Reconstruction loss (how well the model reconstructs the input)
  \item
    KL divergence (a regularization term)
  \item
    Hamiltonian loss (related to the metric tensor)
  \end{itemize}
\item
  \textbf{Optimization}: We use the Adam optimizer with a learning rate
  of 0.001.
\item
  \textbf{Checkpointing}: We save the model state after each epoch,
  allowing us to resume training if needed.
\end{enumerate}

The loss includes several components weighted by the
\texttt{logp\_prefactor} and \texttt{logq\_prefactor} parameters, which
balance reconstruction quality against regularization.

\subsection{Analyzing Training
Results}\label{analyzing-training-results}

After training, we can analyze the training results to understand how
well our model is learning. First, let's load the loss and mean squared
error (MSE) training curves into a dataframe.

\begin{Shaded}
\begin{Highlighting}[]
\CommentTok{\# Loading trained model...}

\CommentTok{\# Find model file}
\NormalTok{model\_file }\OperatorTok{=} \FunctionTok{first}\NormalTok{(Glob.}\FunctionTok{glob}\NormalTok{(}\StringTok{"}\SpecialCharTok{$}\NormalTok{(out\_dir)}\StringTok{/rhvae\_model*.jld2"}\NormalTok{))}
\CommentTok{\# List epoch parameters}
\NormalTok{model\_states }\OperatorTok{=}\NormalTok{ Glob.}\FunctionTok{glob}\NormalTok{(}\StringTok{"}\SpecialCharTok{$}\NormalTok{(state\_dir)}\StringTok{/*.jld2"}\NormalTok{)}

\CommentTok{\# Initialize dataframe to store files metadata}
\NormalTok{df\_meta }\OperatorTok{=}\NormalTok{ DF.}\FunctionTok{DataFrame}\NormalTok{()}

\CommentTok{\# Loop over files}
\ControlFlowTok{for}\NormalTok{ f }\KeywordTok{in}\NormalTok{ model\_states}
    \CommentTok{\# Extract epoch number from file name using regular expression}
\NormalTok{    epoch }\OperatorTok{=} \FunctionTok{parse}\NormalTok{(}\DataTypeTok{Int}\NormalTok{, }\FunctionTok{match}\NormalTok{(}\StringTok{r"epoch}\CharTok{(}\SpecialCharTok{\textbackslash{}d+}\CharTok{)}\StringTok{"}\NormalTok{, f).captures[}\FloatTok{1}\NormalTok{])}
    \CommentTok{\# Load model\_state file}
\NormalTok{    f\_load }\OperatorTok{=}\NormalTok{ JLD2.}\FunctionTok{load}\NormalTok{(f)}
    \CommentTok{\# Extract values}
\NormalTok{    loss\_train }\OperatorTok{=}\NormalTok{ f\_load[}\StringTok{"loss\_train"}\NormalTok{]}
\NormalTok{    loss\_val }\OperatorTok{=}\NormalTok{ f\_load[}\StringTok{"loss\_val"}\NormalTok{]}
\NormalTok{    mse\_train }\OperatorTok{=}\NormalTok{ f\_load[}\StringTok{"mse\_train"}\NormalTok{]}
\NormalTok{    mse\_val }\OperatorTok{=}\NormalTok{ f\_load[}\StringTok{"mse\_val"}\NormalTok{]}
    \CommentTok{\# Generate temporary dataframe to store metadata}
\NormalTok{    df\_tmp }\OperatorTok{=}\NormalTok{ DF.}\FunctionTok{DataFrame}\NormalTok{(}
        \OperatorTok{:}\NormalTok{epoch }\OperatorTok{=\textgreater{}}\NormalTok{ epoch,}
        \OperatorTok{:}\NormalTok{loss\_train }\OperatorTok{=\textgreater{}}\NormalTok{ loss\_train,}
        \OperatorTok{:}\NormalTok{loss\_val }\OperatorTok{=\textgreater{}}\NormalTok{ loss\_val,}
        \OperatorTok{:}\NormalTok{mse\_train }\OperatorTok{=\textgreater{}}\NormalTok{ mse\_train,}
        \OperatorTok{:}\NormalTok{mse\_val }\OperatorTok{=\textgreater{}}\NormalTok{ mse\_val,}
        \OperatorTok{:}\NormalTok{model\_file }\OperatorTok{=\textgreater{}}\NormalTok{ model\_file,}
        \OperatorTok{:}\NormalTok{model\_state }\OperatorTok{=\textgreater{}}\NormalTok{ f,}
\NormalTok{    )}
    \CommentTok{\# Append temporary dataframe to main dataframe}
    \KeywordTok{global}\NormalTok{ df\_meta }\OperatorTok{=}\NormalTok{ DF.}\FunctionTok{vcat}\NormalTok{(df\_meta, df\_tmp)}
\ControlFlowTok{end} \CommentTok{\# for f in model\_states}
\end{Highlighting}
\end{Shaded}

Now, we can plot the training loss and mean squared error (MSE) training
curves.

\begin{Shaded}
\begin{Highlighting}[]
\CommentTok{\# Plotting training loss...}

\CommentTok{\# Initialize figure}
\NormalTok{fig }\OperatorTok{=} \FunctionTok{Figure}\NormalTok{(size}\OperatorTok{=}\NormalTok{(}\FloatTok{600}\NormalTok{, }\FloatTok{300}\NormalTok{))}

\CommentTok{\# Add axis}
\NormalTok{ax }\OperatorTok{=} \FunctionTok{Axis}\NormalTok{(}
\NormalTok{    fig[}\FloatTok{1}\NormalTok{, }\FloatTok{1}\NormalTok{],}
\NormalTok{    xlabel}\OperatorTok{=}\StringTok{"epoch"}\NormalTok{,}
\NormalTok{    ylabel}\OperatorTok{=}\StringTok{"loss"}\NormalTok{,}
\NormalTok{)}

\CommentTok{\# Plot training loss}
\FunctionTok{lines!}\NormalTok{(}
\NormalTok{    ax,}
\NormalTok{    df\_meta.epoch,}
\NormalTok{    df\_meta.loss\_train,}
\NormalTok{    label}\OperatorTok{=}\StringTok{"train"}\NormalTok{,}
\NormalTok{)}
\CommentTok{\# Plot validation loss}
\FunctionTok{lines!}\NormalTok{(}
\NormalTok{    ax,}
\NormalTok{    df\_meta.epoch,}
\NormalTok{    df\_meta.loss\_val,}
\NormalTok{    label}\OperatorTok{=}\StringTok{"validation"}\NormalTok{,}
\NormalTok{)}

\CommentTok{\# Add legend}
\FunctionTok{axislegend}\NormalTok{(ax, position}\OperatorTok{=:}\NormalTok{rt)}

\CommentTok{\# Add axis}
\NormalTok{ax }\OperatorTok{=} \FunctionTok{Axis}\NormalTok{(}
\NormalTok{    fig[}\FloatTok{1}\NormalTok{, }\FloatTok{2}\NormalTok{],}
\NormalTok{    xlabel}\OperatorTok{=}\StringTok{"epoch"}\NormalTok{,}
\NormalTok{    ylabel}\OperatorTok{=}\StringTok{"mean squared error"}\NormalTok{,}
\NormalTok{)}

\CommentTok{\# Plot training loss}
\FunctionTok{lines!}\NormalTok{(}
\NormalTok{    ax,}
\NormalTok{    df\_meta.epoch,}
\NormalTok{    df\_meta.mse\_train,}
\NormalTok{    label}\OperatorTok{=}\StringTok{"train"}\NormalTok{,}
\NormalTok{)}
\CommentTok{\# Plot validation loss}
\FunctionTok{lines!}\NormalTok{(}
\NormalTok{    ax,}
\NormalTok{    df\_meta.epoch,}
\NormalTok{    df\_meta.mse\_val,}
\NormalTok{    label}\OperatorTok{=}\StringTok{"validation"}\NormalTok{,}
\NormalTok{)}

\CommentTok{\# Add legend}
\FunctionTok{axislegend}\NormalTok{(ax, position}\OperatorTok{=:}\NormalTok{rt)}

\NormalTok{fig}
\end{Highlighting}
\end{Shaded}

\includegraphics[width=12.5in,height=6.25in]{rhvae_sim_files/figure-pdf/cell-13-output-1.png}

This code loads all the training checkpoints and plots the training
curves, showing: 1. The total loss over time for both training and
validation sets 2. The Mean Squared Error (MSE) over time for both sets

These plots help us assess if the model is learning effectively and if
it's overfitting or underfitting the data. Since both curves perfectly
track each other, we can conclude that the model is learning effectively
without overfitting the training data.

\subsection{Mapping Data to Latent
Space}\label{mapping-data-to-latent-space}

Now that we have a trained model, we can use it to map our fitness
profiles to the latent space. Again, we will take advantage of the
\texttt{DimArray} type to preserve the metadata of the fitness profiles
when mapping them to the latent space.

\begin{Shaded}
\begin{Highlighting}[]
\CommentTok{\# Mapping data to latent space...}

\CommentTok{\# Standardize the data to have mean 0 and standard deviation 1}
\NormalTok{log\_fitnotype\_std }\OperatorTok{=}\NormalTok{ DD.}\FunctionTok{DimArray}\NormalTok{(}
    \FunctionTok{mapslices}\NormalTok{(slice }\OperatorTok{{-}\textgreater{}}\NormalTok{ StatsBase.}\FunctionTok{transform}\NormalTok{(dt, slice),}
        \FunctionTok{log}\NormalTok{.(fitnotype\_profiles.fitness.data),}
\NormalTok{        dims}\OperatorTok{=}\NormalTok{[}\FloatTok{5}\NormalTok{]),}
\NormalTok{    fitnotype\_profiles.fitness.dims,}
\NormalTok{)}

\CommentTok{\# Load model}
\NormalTok{rhvae }\OperatorTok{=}\NormalTok{ JLD2.}\FunctionTok{load}\NormalTok{(model\_file)[}\StringTok{"model"}\NormalTok{]}
\CommentTok{\# Load latest model state}
\NormalTok{Flux.}\FunctionTok{loadmodel!}\NormalTok{(rhvae, JLD2.}\FunctionTok{load}\NormalTok{(df\_meta.model\_state[}\KeywordTok{end}\NormalTok{])[}\StringTok{"model\_state"}\NormalTok{])}
\CommentTok{\# Update metric parameters}
\NormalTok{AET.RHVAEs.}\FunctionTok{update\_metric!}\NormalTok{(rhvae)}

\CommentTok{\# Define latent space dimensions}
\NormalTok{latent }\OperatorTok{=}\NormalTok{ DD.}\FunctionTok{Dim}\DataTypeTok{\{:latent\}}\NormalTok{([}\OperatorTok{:}\NormalTok{latent1, }\OperatorTok{:}\NormalTok{latent2])}

\CommentTok{\# Map data to latent space}
\NormalTok{dd\_latent }\OperatorTok{=}\NormalTok{ DD.}\FunctionTok{DimArray}\NormalTok{(}
    \FunctionTok{dropdims}\NormalTok{(}
        \FunctionTok{mapslices}\NormalTok{(slice }\OperatorTok{{-}\textgreater{}}\NormalTok{ rhvae.vae.}\FunctionTok{encoder}\NormalTok{(slice).μ,}
\NormalTok{            log\_fitnotype\_std.data,}
\NormalTok{            dims}\OperatorTok{=}\NormalTok{[}\FloatTok{5}\NormalTok{]);}
\NormalTok{        dims}\OperatorTok{=}\FloatTok{1}
\NormalTok{    ),}
\NormalTok{    (log\_fitnotype\_std.dims[}\FloatTok{2}\OperatorTok{:}\FloatTok{4}\NormalTok{]}\OperatorTok{...}\NormalTok{, latent, log\_fitnotype\_std.dims[}\FloatTok{6}\NormalTok{]),}
\NormalTok{)}
\end{Highlighting}
\end{Shaded}

This process involves: 1. Loading the fitness data again and applying
the same preprocessing steps 2. Loading our trained RHVAE model 3. Using
the encoder part of the RHVAE to map each fitness profile to a point in
the 2D latent space 4. Storing the results in a dimensional array that
preserves the metadata (lineage, time, etc.)

We now have a low-dimensional representation of our fitness profiles
that we can visualize and analyze.

\subsection{Visualizing the Latent
Space}\label{visualizing-the-latent-space}

Let's visualize the latent space to see how our model has organized the
fitness profiles:

\begin{Shaded}
\begin{Highlighting}[]
\CommentTok{\# Plotting latent space coordinates...}

\CommentTok{\# Initialize figure}
\NormalTok{fig }\OperatorTok{=} \FunctionTok{Figure}\NormalTok{(size}\OperatorTok{=}\NormalTok{(}\FloatTok{300}\NormalTok{, }\FloatTok{300}\NormalTok{))}

\CommentTok{\# Add axis}
\NormalTok{ax }\OperatorTok{=} \FunctionTok{Axis}\NormalTok{(}
\NormalTok{    fig[}\FloatTok{1}\NormalTok{, }\FloatTok{1}\NormalTok{],}
\NormalTok{    xlabel}\OperatorTok{=}\StringTok{"latent dimension 1"}\NormalTok{,}
\NormalTok{    ylabel}\OperatorTok{=}\StringTok{"latent dimension 2"}\NormalTok{,}
\NormalTok{    aspect}\OperatorTok{=}\FunctionTok{AxisAspect}\NormalTok{(}\FloatTok{1}\NormalTok{)}
\NormalTok{)}

\CommentTok{\# Plot latent space}
\FunctionTok{scatter!}\NormalTok{(}
\NormalTok{    ax,}
    \FunctionTok{vec}\NormalTok{(dd\_latent[latent}\OperatorTok{=}\NormalTok{DD.}\FunctionTok{At}\NormalTok{(}\OperatorTok{:}\NormalTok{latent1)]),}
    \FunctionTok{vec}\NormalTok{(dd\_latent[latent}\OperatorTok{=}\NormalTok{DD.}\FunctionTok{At}\NormalTok{(}\OperatorTok{:}\NormalTok{latent2)]),}
\NormalTok{    markersize}\OperatorTok{=}\FloatTok{5}\NormalTok{,}
\NormalTok{)}

\NormalTok{fig}
\end{Highlighting}
\end{Shaded}

\includegraphics[width=6.25in,height=6.25in]{rhvae_sim_files/figure-pdf/cell-15-output-1.png}

This basic plot shows all points in the latent space. Each point
represents the encoding of a fitness profile. Next, let's colorize the
points by lineage to gain more insight:

\begin{Shaded}
\begin{Highlighting}[]
\CommentTok{\# Plotting latent space coordinates colored by lineage...}

\CommentTok{\# Initialize figure}
\NormalTok{fig }\OperatorTok{=} \FunctionTok{Figure}\NormalTok{(size}\OperatorTok{=}\NormalTok{(}\FloatTok{300}\NormalTok{, }\FloatTok{300}\NormalTok{))}

\CommentTok{\# Add axis}
\NormalTok{ax }\OperatorTok{=} \FunctionTok{Axis}\NormalTok{(}
\NormalTok{    fig[}\FloatTok{1}\NormalTok{, }\FloatTok{1}\NormalTok{],}
\NormalTok{    xlabel}\OperatorTok{=}\StringTok{"latent dimension 1"}\NormalTok{,}
\NormalTok{    ylabel}\OperatorTok{=}\StringTok{"latent dimension 2"}\NormalTok{,}
\NormalTok{    aspect}\OperatorTok{=}\FunctionTok{AxisAspect}\NormalTok{(}\FloatTok{1}\NormalTok{)}
\NormalTok{)}

\CommentTok{\# Loop over lineages}
\ControlFlowTok{for}\NormalTok{ (i, lin) }\KeywordTok{in} \FunctionTok{enumerate}\NormalTok{(DD.}\FunctionTok{dims}\NormalTok{(dd\_latent, }\OperatorTok{:}\NormalTok{lineage))}
    \CommentTok{\# Plot latent space}
    \FunctionTok{scatter!}\NormalTok{(}
\NormalTok{        ax,}
        \FunctionTok{vec}\NormalTok{(dd\_latent[latent}\OperatorTok{=}\NormalTok{DD.}\FunctionTok{At}\NormalTok{(}\OperatorTok{:}\NormalTok{latent1), lineage}\OperatorTok{=}\NormalTok{lin]),}
        \FunctionTok{vec}\NormalTok{(dd\_latent[latent}\OperatorTok{=}\NormalTok{DD.}\FunctionTok{At}\NormalTok{(}\OperatorTok{:}\NormalTok{latent2), lineage}\OperatorTok{=}\NormalTok{lin]),}
\NormalTok{        markersize}\OperatorTok{=}\FloatTok{5}\NormalTok{,}
\NormalTok{        color}\OperatorTok{=}\NormalTok{(ColorSchemes.seaborn\_colorblind[i], }\FloatTok{0.25}\NormalTok{),}
\NormalTok{    )}
\ControlFlowTok{end} \CommentTok{\# for }

\NormalTok{fig}
\end{Highlighting}
\end{Shaded}

\includegraphics[width=6.25in,height=6.25in]{rhvae_sim_files/figure-pdf/cell-16-output-1.png}

This plot shows the latent space colored by lineage. Each lineage is
assigned a distinct color, allowing us to see how different lineages are
distributed in the latent space.

\subsubsection{Visualizing the latent space
curvature}\label{visualizing-the-latent-space-curvature}

One of the main advantages of the RHVAE is that it learns a Riemannian
metric on the latent space, which allows us to relate distances in the
latent space with distances in the original data space, despite the
nonlinear transformations.

The metric tensor, \(\underline{\underline{G}}(\underline{z})\), is a
position-dependent positive-definite matrix that we can evaluate at any
point \(\underline{z}\) in the latent space. For at 2D latent space,
this results in a 2x2 matrix. A way to visualize the local curvature of
the latent space is to compute the so-called metric volume, defined as
the square root of the determinant of the metric tensor,
\(\sqrt{\det(\underline{\underline{G}}(\underline{z}))}\).
\texttt{AutoEncoderToolkit.jl} provides a function to compute the
inverse metric tensor, \(G^{-1}(\underline{z})\), from the model. We can
use this directly to compute the metric volume at each point in the
latent space.

Let's define the ranges of the latent space dimensions and evaluate the
metric tensor at each point in the latent space.

\begin{Shaded}
\begin{Highlighting}[]
\CommentTok{\# Computing metric tensor...}

\CommentTok{\# Define number of points per axis}
\NormalTok{n\_points }\OperatorTok{=} \FloatTok{250}

\CommentTok{\# Extract latent space ranges}
\NormalTok{latent1\_range }\OperatorTok{=} \FunctionTok{range}\NormalTok{(}
    \FunctionTok{minimum}\NormalTok{(dd\_latent[latent}\OperatorTok{=}\NormalTok{DD.}\FunctionTok{At}\NormalTok{(}\OperatorTok{:}\NormalTok{latent1)]) }\OperatorTok{{-}} \FloatTok{1.5}\NormalTok{,}
    \FunctionTok{maximum}\NormalTok{(dd\_latent[latent}\OperatorTok{=}\NormalTok{DD.}\FunctionTok{At}\NormalTok{(}\OperatorTok{:}\NormalTok{latent1)]) }\OperatorTok{+} \FloatTok{1.5}\NormalTok{,}
\NormalTok{    length}\OperatorTok{=}\NormalTok{n\_points}
\NormalTok{)}
\NormalTok{latent2\_range }\OperatorTok{=} \FunctionTok{range}\NormalTok{(}
    \FunctionTok{minimum}\NormalTok{(dd\_latent[latent}\OperatorTok{=}\NormalTok{DD.}\FunctionTok{At}\NormalTok{(}\OperatorTok{:}\NormalTok{latent2)]) }\OperatorTok{{-}} \FloatTok{1.5}\NormalTok{,}
    \FunctionTok{maximum}\NormalTok{(dd\_latent[latent}\OperatorTok{=}\NormalTok{DD.}\FunctionTok{At}\NormalTok{(}\OperatorTok{:}\NormalTok{latent2)]) }\OperatorTok{+} \FloatTok{1.5}\NormalTok{,}
\NormalTok{    length}\OperatorTok{=}\NormalTok{n\_points}
\NormalTok{)}
\CommentTok{\# Define latent points to evaluate}
\NormalTok{z\_mat }\OperatorTok{=} \FunctionTok{reduce}\NormalTok{(hcat, [[x, y] for x }\KeywordTok{in}\NormalTok{ latent1\_range, y }\KeywordTok{in}\NormalTok{ latent2\_range])}

\CommentTok{\# Compute inverse metric tensor}
\NormalTok{Ginv }\OperatorTok{=}\NormalTok{ AET.RHVAEs.}\FunctionTok{G\_inv}\NormalTok{(z\_mat, rhvae)}

\CommentTok{\# Compute metric }
\NormalTok{logdetG }\OperatorTok{=} \FunctionTok{reshape}\NormalTok{(}
    \OperatorTok{{-}}\FloatTok{1} \OperatorTok{/} \FloatTok{2} \OperatorTok{*}\NormalTok{ AET.utils.}\FunctionTok{slogdet}\NormalTok{(Ginv), n\_points, n\_points}
\NormalTok{)}
\end{Highlighting}
\end{Shaded}

Now, let's plot the heatmap of the metric volume and the surface of the
latent space.

\begin{Shaded}
\begin{Highlighting}[]
\CommentTok{\# Plotting latent space metric...}

\CommentTok{\# Initialize figure}
\NormalTok{fig }\OperatorTok{=} \FunctionTok{Figure}\NormalTok{(size}\OperatorTok{=}\NormalTok{(}\FloatTok{700}\NormalTok{, }\FloatTok{300}\NormalTok{))}

\CommentTok{\# Add axis}
\NormalTok{ax1 }\OperatorTok{=} \FunctionTok{Axis}\NormalTok{(}
\NormalTok{    fig[}\FloatTok{1}\NormalTok{, }\FloatTok{1}\NormalTok{],}
\NormalTok{    xlabel}\OperatorTok{=}\StringTok{"latent dimension 1"}\NormalTok{,}
\NormalTok{    ylabel}\OperatorTok{=}\StringTok{"latent dimension 2"}\NormalTok{,}
\NormalTok{    aspect}\OperatorTok{=}\FunctionTok{AxisAspect}\NormalTok{(}\FloatTok{1}\NormalTok{)}
\NormalTok{)}
\NormalTok{ax2 }\OperatorTok{=} \FunctionTok{Axis}\NormalTok{(}
\NormalTok{    fig[}\FloatTok{1}\NormalTok{, }\FloatTok{2}\NormalTok{],}
\NormalTok{    xlabel}\OperatorTok{=}\StringTok{"latent dimension 1"}\NormalTok{,}
\NormalTok{    ylabel}\OperatorTok{=}\StringTok{"latent dimension 2"}\NormalTok{,}
\NormalTok{    aspect}\OperatorTok{=}\FunctionTok{AxisAspect}\NormalTok{(}\FloatTok{1}\NormalTok{)}
\NormalTok{)}

\CommentTok{\# Plot heatmpat of log determinant of metric tensor}
\NormalTok{hm }\OperatorTok{=} \FunctionTok{heatmap!}\NormalTok{(}
\NormalTok{    ax1,}
\NormalTok{    latent1\_range,}
\NormalTok{    latent2\_range,}
\NormalTok{    logdetG,}
\NormalTok{    colormap}\OperatorTok{=}\FunctionTok{Reverse}\NormalTok{(}\FunctionTok{to\_colormap}\NormalTok{(ColorSchemes.PuBu)),}
\NormalTok{)}

\FunctionTok{surface!}\NormalTok{(}
\NormalTok{    ax2,}
\NormalTok{    latent1\_range,}
\NormalTok{    latent2\_range,}
\NormalTok{    logdetG,}
\NormalTok{    colormap}\OperatorTok{=}\FunctionTok{Reverse}\NormalTok{(}\FunctionTok{to\_colormap}\NormalTok{(ColorSchemes.PuBu)),}
\NormalTok{    shading}\OperatorTok{=}\NormalTok{NoShading,}
\NormalTok{    rasterize}\OperatorTok{=}\ConstantTok{true}\NormalTok{,}
\NormalTok{)}

\CommentTok{\# Plot latent space}
\FunctionTok{scatter!}\NormalTok{(}
\NormalTok{    ax2,}
    \FunctionTok{vec}\NormalTok{(dd\_latent[latent}\OperatorTok{=}\NormalTok{DD.}\FunctionTok{At}\NormalTok{(}\OperatorTok{:}\NormalTok{latent1)]),}
    \FunctionTok{vec}\NormalTok{(dd\_latent[latent}\OperatorTok{=}\NormalTok{DD.}\FunctionTok{At}\NormalTok{(}\OperatorTok{:}\NormalTok{latent2)]),}
\NormalTok{    markersize}\OperatorTok{=}\FloatTok{4}\NormalTok{,}
\NormalTok{    color}\OperatorTok{=}\NormalTok{(}\OperatorTok{:}\NormalTok{white, }\FloatTok{0.3}\NormalTok{),}
\NormalTok{    rasterize}\OperatorTok{=}\ConstantTok{true}\NormalTok{,}
\NormalTok{)}

\CommentTok{\# Find axis limits from minimum and maximum of latent points}
\FunctionTok{xlims!}\NormalTok{.(}
\NormalTok{    [ax2, ax1],}
    \FunctionTok{minimum}\NormalTok{(dd\_latent[latent}\OperatorTok{=}\NormalTok{DD.}\FunctionTok{At}\NormalTok{(}\OperatorTok{:}\NormalTok{latent1)]) }\OperatorTok{{-}} \FloatTok{1.5}\NormalTok{,}
    \FunctionTok{maximum}\NormalTok{(dd\_latent[latent}\OperatorTok{=}\NormalTok{DD.}\FunctionTok{At}\NormalTok{(}\OperatorTok{:}\NormalTok{latent1)]) }\OperatorTok{+} \FloatTok{1.5}
\NormalTok{)}
\FunctionTok{ylims!}\NormalTok{.(}
\NormalTok{    [ax2, ax1],}
    \FunctionTok{minimum}\NormalTok{(dd\_latent[latent}\OperatorTok{=}\NormalTok{DD.}\FunctionTok{At}\NormalTok{(}\OperatorTok{:}\NormalTok{latent2)]) }\OperatorTok{{-}} \FloatTok{1.5}\NormalTok{,}
    \FunctionTok{maximum}\NormalTok{(dd\_latent[latent}\OperatorTok{=}\NormalTok{DD.}\FunctionTok{At}\NormalTok{(}\OperatorTok{:}\NormalTok{latent2)]) }\OperatorTok{+} \FloatTok{1.5}
\NormalTok{)}

\CommentTok{\# Add colorbar}
\FunctionTok{Colorbar}\NormalTok{(fig[}\FloatTok{1}\NormalTok{, }\FloatTok{3}\NormalTok{], hm, label}\OperatorTok{=}\StringTok{"√log[det(G̲̲)]"}\NormalTok{)}

\NormalTok{fig}
\end{Highlighting}
\end{Shaded}

\includegraphics[width=14.58333in,height=6.25in]{rhvae_sim_files/figure-pdf/cell-18-output-1.png}

This plot shows the heatmap of the metric volume and the surface of the
latent space. The colorbar shows the metric volume, which is a measure
of the local curvature of the latent space. The darker the region, the
flatter the latent space. One neat feature of the RHVAE is that the
learned metric tensor ``cages'' the data cloud, clearly defining the
regions in which the model is able to make meaningful predictions (since
neural networks are famously terrible at extrapolation). We also see
that within the data cloud, regions of high curvature coincide with the
four regions of low genotype-phenotype fitness we originally defined in
the \href{evolutionary_dynamics.qmd}{Evolutionary Dynamics}, suggesting
that the model is able to capture the underlying phenotypic space.

\subsection{Conclusion}\label{conclusion}

In this notebook, we've trained a Riemannian Hamiltonian Variational
Autoencoder (RHVAE) on synthetic data generated from evolutionary
simulations. We've used the RHVAE to learn a low-dimensional
representation of the fitness profiles and visualized the latent space
to gain insights into the model's organization of the fitness profiles.



\end{document}
